% Aula 04 --- a ponte 02<->04: regularizar ataca variância, não paramétrico ataca viés.
Durante o curso, apresentamos ferramentas que, à primeira vista, andam em sentidos
opostos. A regularização \emph{restringe} o modelo: o Ridge, por exemplo, soma ao
RSS do MQO uma penalidade sobre o tamanho dos coeficientes,
\[
  \widehat{\bbeta}^{\,\text{ridge}}
   = \argmin_{\bbeta}\Big\{\norm{\bm y-\mathbb{X}\bbeta}^2+\lambda\sum_{j=1}^{p}\beta_j^2\Big\}.
\]
Os métodos não paramétricos o \emph{soltam}: o KNN, por exemplo, não supõe forma
nenhuma para $r$ e prevê pela média das respostas dos $k$ vizinhos mais próximos de
$\x$, cujos índices formam o conjunto $\mathcal N_{\x}$,
\[
  \rhat(\x)=\frac1k\sum_{i\in\mathcal N_{\x}} Y_i .
\]
Parece contradição, e não é: são ferramentas para serem usadas em dois contextos
diferentes.

Para ver quais, a referência é a decomposição do erro de previsão num ponto fixo
$\x_0$. Para um estimador $\rhat$ construído com a amostra $\Dados$,
\[
  \E_{\Dados}\!\big[(Y_0-\rhat(\x_0))^2\big]
  = \underbrace{\big(\E_{\Dados}[\rhat(\x_0)]-r(\x_0)\big)^2}_{\text{viés}^2}
  + \underbrace{\Var_{\Dados}\!\big(\rhat(\x_0)\big)}_{\text{variância}}
  + \underbrace{\sigma^2(\x_0)}_{\text{irredutível}},
\]
em que $Y_0$ é uma resposta nova em $\x_0$, $\sigma^2(\x_0)=\Var(Y\mid\X=\x_0)$, e
$\E_{\Dados}$ e $\Var_{\Dados}$ percorrem as amostras de treino possíveis. O último
termo não depende do método; a disputa é entre os dois primeiros.

\begin{itens}
  \item \pts{0,8} Explique, em palavras, o que o viés e a variância medem. Depois
        diga qual dos dois termos a regularização ataca e qual os métodos não
        paramétricos atacam, e o que cada família aceita pagar no outro.
  \item \pts{0,8} Descreva, em uma frase cada, a situação que motiva cada família:
        o que está errado com o modelo \emph{antes} de aplicarmos Ridge, e o que
        está errado \emph{antes} de aplicarmos um spline?
  \item \pts{1,3} O KNN com $k=1$ e o MQO com $p$ grande erram por motivos que,
        na decomposição, caem no mesmo lugar. Explique. Em seguida diga qual é o
        remédio de cada um e por que os dois remédios se parecem.

        Dica para o KNN: trate as covariáveis da amostra como fixas --- assim os
        vizinhos de $\x_0$ também ficam fixos --- e suponha que $Y$ tenha variância
        constante $\sigma^2$ em torno de $r$. Nesse caso a decomposição fica
        \[
          \E\big[(Y_0-\rhat(\x_0))^2\big]
          = \Big(\frac1k\sum_{i\in\mathcal N_{\x_0}} r(\X_i)-r(\x_0)\Big)^2
          + \frac{\sigma^2}{k} + \sigma^2 .
        \]
        O que acontece com cada termo quando $k$ cresce?
  \item \pts{1,1} As duas ideias podem ser combinadas? Dê um exemplo concreto,
        visto no curso, de método que seja ao mesmo tempo não paramétrico e
        regularizado, e diga qual botão controla cada coisa nele.
\end{itens}

\begin{solucao}
\textbf{(a)} Fixe $\x_0$ e pense em $\rhat(\x_0)$ como uma variável aleatória: ele muda
com a amostra $\Dados$.

O \textbf{viés}, $\E_\Dados[\rhat(\x_0)]-r(\x_0)$, compara a \emph{média} das
estimativas, sobre todas as amostras de treino possíveis, com o alvo. É o erro
sistemático, o que sobra depois que o acaso da amostra se cancela na média --- e,
quando a classe de funções do método não contém $r$, ele não some com mais dados. A
\textbf{variância}, $\Var_\Dados(\rhat(\x_0))$, mede a instabilidade: o quanto a
estimativa mudaria se a amostra de treino fosse outra, do mesmo tamanho.

A \textbf{regularização ataca a variância}: ela restringe o espaço de soluções --- no
Ridge, encolhendo os coeficientes --- e aceita pagar viés por isso. Os \textbf{métodos
não paramétricos atacam o viés}: ampliam a classe de funções representáveis, até que
$r$ caiba nela, e aceitam pagar variância. O terceiro termo, $\sigma^2(\x_0)$, fica
fora da disputa: é o erro irredutível do Exercício~\ref{ex:01-ropt}, e nenhuma das
duas famílias mexe nele.

\smallskip
\textbf{(b)} Antes do Ridge, o modelo é \textbf{flexível demais para a amostra}:
covariáveis demais, ou colineares, para o $n$ disponível. A matriz
$\Var(\widehat\bbeta)=\sigma^2(\mathbb{X}^\top\mathbb{X})^{-1}$ tem entradas enormes, e
os coeficientes mudam muito de uma amostra para outra. Excesso de variância.

Antes do spline, o modelo é \textbf{rígido demais para a verdade}: $r$ é curva, e a
reta não a acompanha. O erro que sobra é o erro de aproximação
$\E[(r(\X)-\bbeta^{\ast\top}\X)^2]$ do Exercício~\ref{ex:02-oraculo}, que não depende de
$n$ --- nenhum volume de dados conserta uma reta. Excesso de viés.

"Para a amostra" contra "para a verdade" é o coração da distinção: uma é falta de
dados, a outra é falta de classe.

\smallskip
\textbf{(c)} Os dois erram por \textbf{variância}. Vale ver as contas, porque elas
mostram também por que os remédios se parecem.

\emph{O KNN.} Na situação da dica, os vizinhos de $\x_0$ estão fixos e as respostas são
$Y_i=r(\X_i)+\varepsilon_i$, com ruídos independentes, de média zero e variância
$\sigma^2$. Então $\E[\rhat(\x_0)]=\frac1k\sum_{i\in\mathcal N_{\x_0}}r(\X_i)$ e
$\Var(\rhat(\x_0))=\frac1{k^2}\cdot k\sigma^2=\sigma^2/k$; como $Y_0$ é independente da
amostra, a decomposição da dica sai direto. Com $k=1$, a variância vale $\sigma^2$ ---
tanto quanto o próprio ruído ---, enquanto o viés,
$\big(r(\X_{(1)})-r(\x_0)\big)^2$, é pequeno se $r$ é contínua e o vizinho mais próximo,
$\X_{(1)}$, está perto de $\x_0$. Nos pontos de treino, o ajuste reproduz cada $Y_i$:
resíduo nulo e nenhuma estabilidade.

\emph{O MQO.} No mesmo regime (covariáveis fixas, ruído de variância $\sigma^2$),
$\widehat{\bm y}=\bm H\bm y$ tem matriz de covariância $\sigma^2\bm H\bm H^\top=\sigma^2\bm H$,
porque $\bm H$ é simétrica e idempotente. A variância média dos valores ajustados é
\[
  \frac1n\sum_{i=1}^n\Var\big(\widehat g(\X_i)\big)
  =\frac{\sigma^2}{n}\operatorname{tr}(\bm H)=\sigma^2\,\frac{p+1}{n}.
\]
Ela cresce linearmente com o número de parâmetros e chega a $\sigma^2$ quando
$p+1=n$ --- a interpolação, exatamente como no KNN com $k=1$. Com $p+1>n$, a solução
nem é única.

\emph{Por que caem no mesmo lugar.} O KNN também é um suavizador linear: cada linha da
sua matriz $\bm H$ tem $k$ entradas iguais a $1/k$, de modo que
$\operatorname{tr}(\bm H)=n/k$ (Lista Teórica~04) e
$\Var(\rhat(\X_i))=\sigma^2\cdot k\cdot(1/k)^2=\sigma^2\operatorname{tr}(\bm H)/n$ --- a
mesma expressão do MQO. Nos dois métodos, quem governa a variância é o \emph{número
efetivo de parâmetros} $\operatorname{tr}(\bm H)$: $n/k$ num, $p+1$ no outro. E $k=1$ e
$p+1=n$ são o mesmo desastre --- um parâmetro por observação.

\emph{Os remédios.} O do KNN é \textbf{aumentar $k$}: pela dica, a variância cai como
$\sigma^2/k$, e o preço aparece no primeiro termo, porque vizinhos mais distantes
entram na média e $\frac1k\sum r(\X_i)$ se afasta de $r(\x_0)$. O do MQO é
\textbf{regularizar}. No caso ortonormal do Exercício~\ref{ex:02-ortonormal} a troca é
explícita: o Ridge divide cada coeficiente por $1+\lambda$, o que divide sua variância
por $(1+\lambda)^2$ e cria um viés de $-\lambda\beta_j/(1+\lambda)$.

Os dois se parecem porque fazem a mesma coisa: \textbf{reduzem $\operatorname{tr}(\bm H)$},
trocando variância por viés. No KNN, $\operatorname{tr}(\bm H)=n/k$; no Ridge
ortonormal, $\bm H_\lambda=\mathbb{X}(\mathbb{X}^\top\mathbb{X}+\lambda\mathbb{I})^{-1}\mathbb{X}^\top=\mathbb{X}\mathbb{X}^\top/(1+\lambda)$,
e $\operatorname{tr}(\bm H_\lambda)=p/(1+\lambda)$. Nos dois há um botão --- $k$, que é
inteiro, e $\lambda$, que é contínuo ---, e nos dois ele se escolhe por validação
cruzada, nunca no olho.

E os dois botões correm \textbf{no mesmo sentido}: $k$ grande e $\lambda$ grande dão
modelos mais simples, e na ponta de cima dos dois está o modelo constante --- a média
global quando $k=n$, só o intercepto quando $\lambda\to\infty$. O que corre ao contrário
é o grau do polinômio, ou o número de nós de um spline: ali, grande quer dizer
flexível. O teste seguro, quando a memória falha, é perguntar de que lado está o modelo
constante.

\smallskip
\textbf{(d)} Podem, e o curso tem pelo menos dois exemplos.

\emph{A árvore podada} (Aula~05). O crescimento a torna não paramétrica: o número de
folhas não é fixado de antemão, e cresce com $n$. E a poda por custo--complexidade a
regulariza:
\[
  \sum_{m=1}^{\abs{T}}\ \sum_{i:\;\X_i\in R_m}\big(Y_i-\widehat y_{R_m}\big)^2+\alpha\,\abs{T}
\]
é "ajuste mais preço vezes complexidade", a mesma estrutura do Ridge, com o número de
folhas no lugar da soma dos quadrados dos coeficientes. A flexibilidade disponível vem
do crescimento; quanto dela é usado, decide o $\alpha$.

\emph{As smoothing splines} (Aula~04), que põem um nó em cada observação --- não
paramétricas ao extremo, com $n$ parâmetros nominais --- e controlam a flexibilidade
penalizando a curvatura:
\[
  \sum_{i=1}^n\big(y_i-g(x_i)\big)^2+\gamma\int\big(g''(x)\big)^2\,dx .
\]
A base cresce com $n$ (o lado não paramétrico), e $\gamma$ regula quanta daquela
liberdade é de fato usada (o lado regularizado). Os nós deixaram de ser o botão ---
são todos ---, e quem manda é $\gamma$: quando ele vai de $0$ a $\infty$, os graus de
liberdade efetivos caem de $n$ para $2$, os da reta de mínimos quadrados.

(Vale também, no mesmo espírito, aplicar Ridge ou Lasso a uma base grande --- muitos
monômios, ou um spline com muitos nós.)
\end{solucao}
